\documentclass[11pt]{article}

\usepackage[margin=0.9in]{geometry}
\usepackage{amsmath,amssymb,amsthm,mathtools}
\usepackage{graphicx}
\usepackage{microtype}
\usepackage{xcolor}
\usepackage[hidelinks]{hyperref}

\newtheorem{theorem}{Theorem}
\newtheorem{lemma}{Lemma}
\newcommand{\R}{\mathbb R}
\newcommand{\esssup}{\operatorname*{ess\,sup}}
\newcommand{\Lip}{\operatorname{Lip}}
\newcommand{\vol}{\operatorname{vol}}

\title{The Liminf Besov Criterion Characterizes Lipschitz Functions\\
\large A full positive answer to Remark 2.13 of arXiv:1610.05162}
\author{Automated mathematics run \texttt{fa\_banach\_001}, lane 14}
\date{21 July 2026}

\begin{document}
\maketitle

\begin{center}
\fcolorbox{black}{gray!8}{\parbox{0.88\textwidth}{
\textbf{Status.} Candidate full solution, likely valid, pending human review.
The theorem answers the source's liminf question for every dimension
$N\geq1$ and every $1\leq q<\infty$.
}}
\end{center}

\section{Source question}

For $0<r<1$, $1\leq q<\infty$, and $f\in L^\infty(\R^N)$, Brasseur
\cite{Brasseur} uses the first-difference Besov seminorm
\begin{equation}
 [f]_{B^r_{\infty,q}}^q
 =\int_{\R^N}
   \frac{\|f(\cdot+h)-f\|_\infty^q}{|h|^{N+rq}}\,dh.
 \label{eq:besov}
\end{equation}
Theorem 2.12 in the source proves that
\[
 \limsup_{r\uparrow1}(1-r)^{1/q}\|f\|_{B^r_{\infty,q}}<\infty
\]
implies that $f$ has a bounded Lipschitz representative, and compares that
limsup with the Lipschitz seminorm. Remark 2.13 asks whether the limsup in
both assertions can be replaced by a liminf. The complete source statement
appears on p.~9 of the arXiv PDF:

\begin{figure}[htbp]
 \centering
 \includegraphics[width=\textwidth]{figures/open_problem_crop.png}
 \caption{Full-width crop of Theorem 2.12 and Remark 2.13 in the source.}
 \label{fig:source}
\end{figure}

We give a positive answer.

\section{Proof intuition}

Write $\omega(h)=\|f(\cdot+h)-f\|_\infty$. A tempting obstruction is that
the normalized Besov norms might have deep valleys as $r\uparrow1$. The
translation modulus prevents this. Its triangle inequality
$\omega(a)\leq\omega(h)+\omega(a-h)$ spreads a large difference quotient at
one displacement across a whole ball of smaller displacements. Subdividing
that displacement shows that this spreading occurs at every smaller radius.

Consequently, if $\omega(a_0)/|a_0|$ is at least $M$, then the cumulative
mass $F(t)=\int_{|h|\leq t}\omega(h)^q\,dh$ is bounded below by a constant
times $M^q t^{N+q}$ at all sufficiently small $t$. Integrating the Besov
kernel against this cumulative bound produces
$M^q/(1-r)$. Hence the normalized liminf is bounded below by a fixed
multiple of $M$. Letting $M$ increase proves that any non-Lipschitz function
has infinite normalized liminf.

\section{Translation modulus lemmas}

For $f\in L^\infty(\R^N)$ set
\[
 \omega(h)=\|f(\cdot+h)-f\|_\infty,
 \qquad
 L(f)=\sup_{h\ne0}\frac{\omega(h)}{|h|}\in[0,\infty].
\]

\begin{lemma}[Finite translation seminorm]
If $L(f)<\infty$, then the $L^\infty$ class of $f$ has a bounded
$L(f)$-Lipschitz representative. Conversely, a bounded $L$-Lipschitz
representative satisfies $L(f)\leq L$.
\end{lemma}

\begin{proof}
Only the first assertion needs comment. Let $\varphi_\varepsilon$ be a
standard smooth approximate identity and set
$f_\varepsilon=f*\varphi_\varepsilon$. Translation commutes with convolution,
so
\[
 \|f_\varepsilon(\cdot+h)-f_\varepsilon\|_\infty
 \leq\omega(h)\leq L(f)|h|.
\]
Thus the $f_\varepsilon$ are uniformly bounded and equi-Lipschitz. A diagonal
Arzela--Ascoli argument gives a locally uniformly convergent subsequence with
a bounded $L(f)$-Lipschitz limit $g$. On the other hand,
$f_\varepsilon\to f$ in $L^1_{\mathrm{loc}}$, hence $g=f$ almost everywhere.
The converse follows directly from the pointwise Lipschitz inequality.
\end{proof}

Let $v_N=\vol(B(0,1))$ and
\[
 F(t)=\int_{|h|\leq t}\omega(h)^q\,dh.
\]

\begin{lemma}[Cumulative spreading]
Suppose $a_0\ne0$ and
$\omega(a_0)/|a_0|\geq M$. Then, for every $0<s\leq2|a_0|$,
\begin{equation}
 F(s)\geq c_* M^q s^{N+q},
 \qquad c_*=v_N2^{-(N+3q)}.
 \label{eq:cumulative}
\end{equation}
\end{lemma}

\begin{proof}
The translation triangle inequality gives, for all $a,h\in\R^N$,
\[
 \omega(a)\leq\omega(h)+\omega(a-h).
\]
If $|a|\leq t$, raise this inequality to the $q$th power and integrate over
$h\in B(0,t)$. Since $a-B(0,t)\subset B(0,2t)$,
\begin{equation}
 v_Nt^N\omega(a)^q\leq2^qF(2t).
 \label{eq:average}
\end{equation}

Fix $0<t\leq|a_0|$ and put $n=\lceil|a_0|/t\rceil$ and $a=a_0/n$.
Then $|a|\leq t$, $|a|\geq t/2$, and subadditivity along the segment from
$0$ to $a_0$ gives
\[
 \omega(a)\geq\frac{\omega(a_0)}n\geq M|a|\geq\frac{Mt}{2}.
\]
Substitution in \eqref{eq:average}, followed by $s=2t$, yields
\eqref{eq:cumulative}.
\end{proof}

\section{Main theorem}

\begin{theorem}[Liminf replacement]
Let $N\geq1$ and $1\leq q<\infty$. For every
$f\in L^\infty(\R^N)$, with extended-real values allowed,
\begin{equation}
 c_{N,q}L(f)
 \leq\liminf_{r\uparrow1}(1-r)^{1/q}\|f\|_{B^r_{\infty,q}}
 \leq\limsup_{r\uparrow1}(1-r)^{1/q}\|f\|_{B^r_{\infty,q}}
 \leq C_{N,q}L(f),
 \label{eq:main}
\end{equation}
where one may take
\[
 c_{N,q}=\left(v_N2^{-(N+3q)}\frac{N+q}{q}\right)^{1/q},
 \qquad
 C_{N,q}=\left(\frac{Nv_N}{q}\right)^{1/q}.
\]
In particular,
\[
 \liminf_{r\uparrow1}(1-r)^{1/q}\|f\|_{B^r_{\infty,q}}<\infty
 \quad\Longleftrightarrow\quad
 f\text{ has a bounded Lipschitz representative}.
\]
Thus every limsup in Theorem 2.12 and Remark 2.13 of the source can be
replaced by a liminf, up to constants depending only on $N$ and $q$.
\end{theorem}

\begin{proof}
Write $A_r=[f]_{B^r_{\infty,q}}^q$ and fix $a_0\ne0$ with
$\omega(a_0)/|a_0|\geq M$. Put $R=2|a_0|$ and
$\alpha=N+rq$. If $A_r=\infty$, the desired lower estimate at this $r$ is
automatic. Otherwise, Stieltjes integration by parts gives
\begin{align}
 A_r
 &\geq\int_{0}^{R}s^{-\alpha}\,dF(s)\\
 &=R^{-\alpha}F(R)+\alpha\int_0^R F(s)s^{-\alpha-1}\,ds.
 \label{eq:parts}
\end{align}
The boundary term at zero vanishes: indeed,
\[
 s^{-\alpha}F(s)
 \leq\int_{|h|\leq s}|h|^{-\alpha}\omega(h)^q\,dh\longrightarrow0
\]
by absolute continuity of the finite integral. Using
Lemma 2 in \eqref{eq:parts},
\begin{align}
 A_r
 &\geq \alpha c_*M^q
   \int_0^R s^{N+q-\alpha-1}\,ds
 =\frac{\alpha c_*}{q(1-r)}M^qR^{q(1-r)}.
 \label{eq:lower}
\end{align}
It follows that
\[
 \liminf_{r\uparrow1}(1-r)A_r
 \geq c_*\frac{N+q}{q}M^q.
\]
Taking the supremum over all admissible $M$ proves the first inequality in
\eqref{eq:main}; if $L(f)=\infty$, the same argument shows that the liminf is
infinite. The inhomogeneous $L^\infty$ term disappears after multiplication
by $(1-r)^{1/q}$.

For the upper bound assume $L(f)<\infty$ and put $B=\|f\|_\infty$. Then
$\omega(h)\leq L(f)|h|$ and $\omega(h)\leq2B$. Splitting at $|h|=1$ and
using polar coordinates,
\begin{align*}
 A_r
 &\leq Nv_NL(f)^q\int_0^1\rho^{q(1-r)-1}\,d\rho
   +Nv_N(2B)^q\int_1^\infty\rho^{-rq-1}\,d\rho\\
 &=\frac{Nv_N}{q(1-r)}L(f)^q+\frac{Nv_N}{rq}(2B)^q.
\end{align*}
Multiplication by $1-r$ and passage to the limsup proves the last inequality
in \eqref{eq:main}. Lemma 1 identifies finiteness of $L(f)$ with the source's
Lipschitz conclusion.
\end{proof}

\section{Verification, novelty, and limitations}

The proof uses only the first-difference definition \eqref{eq:besov}; no
Littlewood--Paley equivalence with constants depending on $r$ is invoked. The
script \texttt{code/verify\_constants.py} checks the displayed constants and
the normalized radial integral for the model subadditive modulus
$\omega(h)=\min\{|h|,1\}$. These computations are sanity checks and do not
replace the cumulative or integration-by-parts arguments.

The local run indexes, two bounded web-search rounds on 21 July 2026, and
full-text inspection of five cached arXiv papers citing the source used the
exact Remark 2.13 sentence and close variants involving \emph{liminf},
\emph{translation modulus}, \emph{Abel mean}, \emph{Besov norm}, and
\emph{Lipschitz regularity}. No later answer or matching proof was located.
This is not an exhaustive citation search, so novelty confidence is moderate.

The theorem addresses exactly the $p=\infty$, $r\uparrow1$ question in
Remark 2.13. It does not alter the source's separate negative theorem for
fractional target smoothness, nor does it claim an exact limit or optimal
constants. Human review should focus on Lemma 2, the boundary term in
\eqref{eq:parts}, and the representative argument in Lemma 1.

\begin{thebibliography}{9}
\bibitem{Brasseur}
J. Brasseur,
\emph{A Bourgain--Brezis--Mironescu characterization of higher order
Besov--Nikol'skii spaces},
Ann. Inst. Fourier (Grenoble) 68 (2018), 1671--1714;
arXiv:\href{https://arxiv.org/abs/1610.05162}{1610.05162}.
\end{thebibliography}

\end{document}
